\section{Magnetometer (LSM9DS1)}

\subsection{Overview}
The LSM9DS1 magnetometer measures magnetic field strength along three axes, enabling orientation and heading detection. Integrated into the Arduino Nano 33 BLE Sense, it supports applications such as navigation, spatial tracking, and motion sensing \cite{Bosch:2024c, STM:2015}.

\subsection{Technical Specifications}
The magnetometer provides the following key features \cite{STM:2015}:
\begin{itemize}
    \item \textbf{Full-scale range:} $\pm4$, $\pm8$, $\pm12$, $\pm16$ gauss.
    \item \textbf{Sensitivity:} 0.14 mgauss/LSB at $\pm4$ gauss.
    \item \textbf{Output data rate (ODR):} 0.625 Hz to 80 Hz.
    \item \textbf{Power consumption:} 2 mA in high-performance mode.
    \item \textbf{Operating temperature:} $-40^{\circ}$C to $+85^{\circ}$C.
    \item \textbf{Interface:} I\textsuperscript{2}C and SPI compatible.
\end{itemize}

\subsection{Calibration Procedure}
Calibration is essential to account for magnetic distortions:
\begin{enumerate}
    \item Rotate the sensor slowly in all directions to collect data.
    \item Compute offsets and scaling factors for each axis.
    \item Apply calibration values in the firmware to correct readings.
\end{enumerate}

\subsection{Calibration Code}
The following example demonstrates magnetometer calibration:
\begin{lstlisting}[language=C++, caption=Magnetometer Calibration Code]
#include <Arduino_LSM9DS1.h>

float offsetX = 0, offsetY = 0, offsetZ = 0;

void calibrateMagnetometer() {
    const int numSamples = 200;
    float minX = 1e6, maxX = -1e6;
    float minY = 1e6, maxY = -1e6;
    float minZ = 1e6, maxZ = -1e6;

    for (int i = 0; i < numSamples; i++) {
        if (IMU.magneticFieldAvailable()) {
            float mx, my, mz;
            IMU.readMagneticField(mx, my, mz);
            minX = min(minX, mx);
            maxX = max(maxX, mx);
            minY = min(minY, my);
            maxY = max(maxY, my);
            minZ = min(minZ, mz);
            maxZ = max(maxZ, mz);
        }
        delay(50);
    }

    offsetX = (maxX + minX) / 2;
    offsetY = (maxY + minY) / 2;
    offsetZ = (maxZ + minZ) / 2;
}

void setup() {
    Serial.begin(9600);
    while (!Serial);

    if (!IMU.begin()) {
        Serial.println("Failed to initialize magnetometer!");
        while (1);
    }

    Serial.println("Calibrating magnetometer...");
    calibrateMagnetometer();
    Serial.println("Calibration complete.");
}

void loop() {
    float mx, my, mz;

    if (IMU.magneticFieldAvailable()) {
        IMU.readMagneticField(mx, my, mz);
        mx -= offsetX;
        my -= offsetY;
        mz -= offsetZ;

        Serial.print("Calibrated Magnetic Field (X, Y, Z): ");
        Serial.print(mx, 2); Serial.print(" \ensuremath{\mu\text{T}}, ");
        Serial.print(my, 2); Serial.print(" \ensuremath{\mu\text{T}}, ");
        Serial.print(mz, 2);
    }
    delay(100);
}
\end{lstlisting}

\subsection{Simple Code Example}
Here is a basic code example to read magnetic field strength:
\begin{lstlisting}[language=C++, caption=Magnetometer Simple Code]
#include <Arduino_LSM9DS1.h>

void setup() {
    Serial.begin(9600);
    while (!Serial);

    if (!IMU.begin()) {
        Serial.println("Failed to initialize magnetometer!");
        while (1);
    }

    Serial.println("Magnetometer initialized.");
}

void loop() {
    float mx, my, mz;

    if (IMU.magneticFieldAvailable()) {
        IMU.readMagneticField(mx, my, mz);
        Serial.print("Magnetic Field (X, Y, Z): ");
        Serial.print(mx); Serial.print(" \ensuremath{\mu\text{T}}, ");
        Serial.print(my); Serial.print(" \ensuremath{\mu\text{T}}, ");
        Serial.println(mz); Serial.print(" \ensuremath{\mu\text{T}}");
    }
    delay(100);
}
\end{lstlisting}

\subsection{Testing}
To test the magnetometer:
\begin{itemize}
    \item Compare readings with a calibrated reference compass.
    \item Rotate the sensor in different orientations and observe stable outputs.
\end{itemize}

\subsection{Further Readings}
\begin{itemize}
    \item LSM9DS1 Datasheet \cite{STM:2015}.
    \item Arduino Nano 33 BLE Sense Documentation \cite{ArduinoNano33Sense2025}.
    \item Adafruit LSM9DS1 Guide \cite{AdafruitLSM9DS1}.
\end{itemize}